\begin{abstract}
We present \model, a model capable of realistic video synthesis, given a sequence of textual prompts. Generating videos from text is particularly challenging due to the computational cost, limited quantities of high quality text-video data and variable length of videos. To address these issues, we introduce a new model for learning video representation which compresses the video to a small representation of discrete tokens. This tokenizer uses causal attention in time, which allows it to work with variable-length videos. 
To generate video tokens from text we are using a bidirectional masked transformer conditioned on pre-computed text tokens. The generated video tokens are subsequently de-tokenized to create the actual video. To address data issues, we demonstrate how joint training on a large corpus of image-text pairs as well as a smaller number of video-text examples can result in generalization beyond what is available in the video datasets. Compared to the previous video generation methods, \model can generate arbitrary long videos conditioned on a sequence of prompts (i.e. time variable text or \textit{a \story}) in open domain. To the best of our knowledge, this is the first time a paper studies generating videos from time variable prompts. In addition, compared to the per-frame baselines, the proposed video encoder-decoder computes fewer tokens per video but results in better spatio-temporal consistency.


% \pikinder{Sep 20: Me thinking out loud.
% \begin{itemize}
%     \item First effective video approaches NUWA and CogVideo rely on a per frame encoding.
%     \item VQGAN sota for image, but artefacts in video -> CVIVIT
%     \item Video data is limited, therefore need to train on image and text -> CVIVIT
%     \item Video sequences are long -> Generation with transformers slow -> Maskgit, CVIVIT
%     \item Story mode is not in the abstract
% \end{itemize}
% }
\end{abstract}